\section{Implementability and encoding}\label{sec:implementation}

The preceding construction specifies the nuclear configuration, and
Section~\ref{subsec:physical-thresholds} gives threshold accuracies
that preserve the source promise. It remains to compute those outputs
in deterministic polynomial time. We first fix the order of the scales,
then explain how to evaluate the orbital integrals and transported nodes
to the required accuracy.

\subsection{Order and uniformity of the scale choices}

Fix the bounded well data and the integer $\rho$ using the positivity
bound in~\eqref{eq:4}, under the conditions that the well supports are
disjoint and $b_i/\lambda\le1$. As explained there, that derivative
bound uses only $2,W,\zeta$, so it does not depend on the later Gaussian
or Coulomb coefficients. Next fix the source cutoff and three spin
gadget penalty scales from the source size and gap bounds. Fix the
inverse-polynomial rational $\tau$ using the uniform Coulomb coercivity
constant for sufficiently large $D$. Finally choose
\[
 k=\lceil Cn^A\rceil,\qquad h=k^{-1},\qquad
 \lambda=8k^3/\rho,
\]
with sufficiently large absolute positive integers $A,C$.

At this last stage, $K_e,1/K_e$, the layout extent, and $b_i$ already
have fixed polynomial size bounds. Their actual values, including the
centers and $b_i$, may be computed after $k$ is fixed. Each remaining
smallness condition has the form
$p(n,\log\lambda)\lambda^{-\alpha}$ bounded by a specified inverse
polynomial in $n$, with fixed $\alpha>0$. The endpoint brackets for
hopping calibration have the same fixed-power character. A sufficiently
large fixed $A$ therefore meets every asymptotic requirement; increasing
$C$ covers all $n\ge2$ and the fixed separation and gap conditions.
The constants whose existence was proved in the preceding estimates
can be dominated by absolute constants in these choices.

The uniformity of this argument depends on the bounds already proved:
the polynomials $p(n,D)$ come from fixed-order derivatives and moments,
fixed-dimensional wells and gadgets, and polynomially many matrix
entries. The large box is treated separately. Its effects were bounded
in~\eqref{eq:5}, by orbital tails, and by the summable fourth
derivatives in the cubature argument. Its total mass enters only the
much smaller node-rounding error. Thus no exponent needs to grow with
the input, and $H$, the box volume, and $1/\delta_R$ are all polynomial
in $n$ after the fixed choice of $A$.

\subsection{Numerical evaluation and hopping calibration}

All required real quantities can be evaluated to inverse-polynomial
absolute accuracy. This includes the large scalar shift in
\eqref{eq:10}: since $C_H$ is multiplied by $\lambda^2m$, evaluating
it within
\[
             \frac{\tau^2\gamma}{128\lambda m}
\]
uses at most one quarter of the allowed threshold error
$\lambda\tau^2\gamma/32$. The other scalar terms can be allocated
the remaining budget in the same way, because their amplification
factors and their number are polynomially bounded. The magnitude
$|C_H|\le CSH$ therefore causes no demand for exponentially fine
absolute precision.

We give elementary integration procedures for the quantities used in
the construction. In the Gaussian heat formula for $\phi$ and its
derivatives, substitute $t=s^2$ and rescale the Gaussian integration
variable. For a derivative $D^\alpha$ of any required order, the result
is
\[
 D^\alpha\phi(r)=\frac1\pi
 \int_0^\infty\!\int_{\mathbb R^2}
       2s e^{-s^2}e^{-|a|^2}D^\alpha f(r+2sa)\,da\,ds.
\]
This integrand is regular at $s=0$. The needed derivatives of $f$ are
bounded piecewise polynomials with bounded first derivatives. Polynomial
integration cutoffs make the Gaussian tails smaller than any prescribed
inverse-polynomial error; on the remaining domain the integrand and its
Lipschitz bound are polynomially bounded. Elementary grid sums in this
fixed dimension therefore suffice. Division by $\phi$ in $W$ and its
derivatives occurs only on the unit disk, where $\phi$ has a fixed
positive lower bound. The other decaying orbital integrals, including
normalization and hopping integrals, can likewise be truncated about
their centers with polynomial cutoffs.

For Coulomb integrations, including $v_{ij}$ and $C_H$, replace
$|x|^{-1}$ inside radius $\varepsilon$ by $\varepsilon^{-1}$, with
$\varepsilon$ an appropriately small inverse polynomial. Integrating
the local singularity against bounded densities gives an error of order
$\varepsilon^2$ times the relevant polynomial volume bounds. The capped
kernel has Lipschitz constant at most $\varepsilon^{-2}$, so grid sums
again give the required accuracy in polynomial time. All integrations
have fixed dimension and fixed nesting depth. Exponentials, logarithms,
square roots, and the other elementary constants can also be evaluated
by rational approximations to the required accuracy; the arguments and
reciprocal absolute tolerances are polynomially bounded. Square roots
at zero can be handled by bisection. Choosing guard precision for the
nested evaluations keeps their accumulated errors within the allotted
budgets.

The link hopping equations have strict endpoint brackets, from
exponential decay at the two prescribed distance scales and polynomial
upper and inverse lower bounds on their targets. In approximate
bisection, evaluate the residual to a small fraction of its per-link
tolerance. If the sign is ambiguous, the residual is already small
enough. Otherwise retain the sign-change bracket until its width times
the Lipschitz bound is within tolerance. No derivative lower bound is
needed. The fixed layout contractions and subsequent rational center
approximations can be computed more accurately than these tolerances
require.

\subsection{Flow evaluation and rational coordinates}

Once the rational centers $u_i$ are fixed, their exact real coefficients
$b_i$ define a reference density and a reference flow $G$. We approximate
the images under that flow; numerical evaluation of $b_i$ is part of
the velocity evaluation error. Thus all numerical and final rational
rounding errors can be charged to the single displacement budget
$\delta_R$, without changing the reference density or any cell mass.

The flow lasts for time one and has uniformly bounded velocities and
first derivatives in space and time. Euler stepping with an
inverse-polynomial time mesh therefore suffices. More explicitly, for
mesh $\nu$ (the reciprocal of the number of steps), initial error $a_0$,
per-step arithmetic error $a_1$, and velocity evaluation error
$\epsilon_v$, the final position error is
\[
                  O(a_0+\nu+\epsilon_v+a_1/\nu).
\]
Indeed a step multiplies the preceding error by at most $1+C\nu$ and
adds $O(\nu^2+\nu\epsilon_v+a_1)$. Each parameter can be chosen
inverse polynomially small so that the final error is at most
$\delta_R$. Rounding intermediate positions on sufficiently fine
rational grids keeps their bit lengths polynomial. Near support
boundaries, the specified smooth vanishing and derivative bounds give
the same error control. None of these operations uses a many-electron
wavefunction.

\subsection{Output size and encoding}

There are $8|\Omega|h^{-3}$ point nuclei, polynomially many in $n$,
and every nucleus has charge one. Their rational physical positions,
the unary electron number, and the rational thresholds from
Section~\ref{subsec:physical-thresholds} can all be computed with
polynomial encoding lengths. Reducing rational outputs to lowest terms
gives the specified input format. To define the map on every input,
one may first check the source syntax and the fixed polynomial bounds
used in the construction, and output a fixed instance when they fail;
this does not affect either promised implication.

The deterministic construction thus supplies the outputs whose energy
comparison and threshold margins were proved in~\eqref{eq:10} and the
following inequalities. With the electronic-energy convention that
omits nuclear repulsion, this completes the proof of
Theorem~\ref{thm:main}, with $Z_{\max}=1$ and $c=1$.
